\section{Background}
\label{sec:rs-background}

The Recursive Spawning Protocol (RSP) operates at the intersection of several research traditions: agent lifecycle management in multi-agent AI systems, accountability and provenance in distributed computing, fixed versus mutable delegation structures, and hierarchical task decomposition. This section situates RSP within each tradition, traces the conceptual gaps that motivate it, and establishes the BX3 Framework as the minimal sufficient substrate for RSP's accountability guarantees.

% -------------------------------------------------------
\subsection{Agent Lifecycle Management in Multi-Agent Systems}
% -------------------------------------------------------

Multi-agent systems (MAS) have evolved from constrained research testbeds into the dominant operational paradigm for deploying autonomous AI in production environments. This shift carries accountability consequences absent from classical MAS research. Classical agent lifecycle management was primarily an engineering concern: spawning, scheduling, and terminating agents to maximize throughput or minimize latency, with accountability residing with the system operator who configured the agent population \cite{bansal2019}. As AI agents move into healthcare, infrastructure control, and financial decision-making, this assumption becomes untenable. When a spawned agent's action produces real-world consequences, the question of who authorized the spawn, through what chain of delegation, and with what scope constraints is no longer academic --- it is a legal, regulatory, and safety question with no clean answer in conventional architectures.

The literature distinguishes two broad spawning paradigms. The \emph{design-time paradigm}, embodied by early frameworks such as the Agent Modeling Language (AML), treats spawning as a static configuration decision: an agent either could or could not spawn children based on its implemented capabilities, with no runtime enforcement of spawn justification or depth constraints \cite{bansal2019}. The \emph{runtime paradigm}, exemplified by modern agentic frameworks, treats spawning as dynamically discovered: agents spawn children based on perceived task requirements encountered during execution, with the parent chain growing organically as the system encounters novel subproblems \cite{autoGPT2023, langchain2023}. This dynamic extension enables unprecedented task flexibility but introduces a structural accountability gap: the delegation chain is mutable, unrecorded at the protocol level, and potentially unbounded.

Recent work has begun addressing this gap directly. The Human Delegation Provenance (HDP) protocol identifies the core vulnerability in multi-hop delegations: in human $\rightarrow$ orchestrator $\rightarrow$ sub-agent $\rightarrow$ tool-agent chains, verifying that terminal actions were genuinely authorized by a human requires reconstructing the full delegation chain --- a task current standards (OAuth 2.0, JWTs, UCAN) do not adequately address for append-only, human-provenance requirements in agentic systems \cite{hdp-draft, hdp-draft-2}. HDP proposes cryptographic tokens capturing each delegation hop as a signed entry in an append-only chain, enabling offline verification using only the issuer's public key and session ID. The Agent Lifecycle Protocol (ALP) extends this by defining canonical lifecycle events --- Genesis, Fork, Migration, Retirements, Succession, and Decommission --- with explicit state-machine semantics, transition rules, and boundary hooks, along with a genealogical registry capturing both genetic (model, architecture) and epigenetic (config, memory, reputation history) lineage for complete provenance across delegation chains \cite{alp2025}. PROV-AGENT further demonstrates this convergence, introducing a provenance model extending W3C PROV with the Model Context Protocol to link agent decisions and prompts with end-to-end workflow data, enabling near real-time provenance capture and cross-facility evaluation \cite{prov-agent2025}.

Unified Agent Lifecycle Management (UALM) provides a complementary five-layer blueprint for end-to-end governance of agent fleets --- identity and persona registry, orchestration and cross-domain mediation, PHI-bounded context and memory, runtime policy enforcement with kill-switch triggers, and lifecycle management linked to credential revocation and audit logging \cite{ualm2025}. Together, these protocols demonstrate that the industry has identified the accountability gap in agentic lifecycles and is converging on cryptographic, append-only delegation records as the solution. RSP's contribution is providing the structural enforcement layer that makes these delegation records immutable by construction, rather than conventionally immutable by policy.

The core problem RSP addresses is \emph{spawn chain opacity}: the inability to determine, at any point during or after execution, whether a given agent's authority traces to a human principal or to an autonomous predecessor. This opacity arises because parent pointers are mutable in conventional architectures --- a child agent's authority can be redirected after provisioning --- and because logging is typically append-only but not tamper-evident. RSP's immutable parent pointer and Fact Layer write protocol make chain opacity architecturally impossible.

% -------------------------------------------------------
\subsection{Accountability and Provenance in Distributed Systems}
% -------------------------------------------------------

Provenance tracking in distributed systems is a well-studied problem with a well-studied gap: provenance records are typically metadata annotations on data objects, not runtime constraints on system behavior. The W3C PROV model defines a formal vocabulary for recording the origin of data products, the agents responsible for them, and the processes that transformed them \cite{moreau2015-prov}. PROV is powerful for forensic reconstruction but operates on the assumption that provenance records are created voluntarily and can be retroactively modified without architectural consequence --- a compromised actor can alter PROV records to obscure unauthorized actions, with modification undetectable without additional tamper-evidence mechanisms.

Blockchain-based provenance systems address the tamper-evidence gap by making records immutable once written. Hyperledger Fabric provides an append-only ledger where every record is hash-chained to its predecessors, making retroactive insertion, deletion, or modification detectable by any participant with read access \cite{hyperledger2015}. Bitcoin's Nakamoto consensus extends this with a decentralized timestamp server that makes double-spend attacks computationally infeasible \cite{nakamoto2008bitcoin}. These systems demonstrate that tamper-evident provenance is achievable at scale but do not address what \emph{constrains} the actors that write to the chain --- they assume write authorization has already been established and is not itself a target for compromise. The trustless provenance tree framework formalizes this limitation: achieving complete provenance integrity requires three mechanisms working together --- cryptographic priority (binding commitments to tree state), governance cascade (off-chain or on-chain governance), and contract enforcement (on-chain rules and audits). No single mechanism suffices \cite{trustless2026}.

In multi-agent AI systems, the accountability problem has an additional dimension absent from conventional distributed systems: the agent itself may not be a stable, identifiable entity. An LLM-based agent's identity is typically derived from its model weights, prompt configuration, and tool access --- none of which are inherently stable across sessions or updates. Recent work on agent auditability identifies five dimensions critical for accountability: action recoverability, lifecycle coverage, policy checkability, responsibility attribution, and evidence integrity. The analysis demonstrates that no agent system is truly accountable without auditability --- and that auditability requires enforcing provenance as a structural property, not a logging convention \cite{auditable-agents2026}. RSP addresses the identity problem by defining agent identity through the \emph{spawn warrant} --- a cryptographically signed artifact issued by the Fact Layer at provisioning encoding the agent's parent, depth, declared purpose, and activation timestamp. The warrant is the agent's identity credential; without a valid warrant, the agent is not authorized to act.

The distinction between mutable and immutable delegation chains maps to the distinction between \emph{attribution} and \emph{enforcement} in accountability design. Mutable chains enable retrospective attribution: after an event, one can trace backwards to identify responsible parties. Immutable chains enable runtime enforcement: the system can prevent unauthorized actions because the chain of authority is structurally fixed and continuously verifiable. RSP's contribution is demonstrating that immutable delegation chains, when implemented as a Fact Layer constraint rather than a logging convention, enable runtime enforcement of the accountability chain --- making the guarantee structural rather than forensic.

% -------------------------------------------------------
\subsection{Fixed vs. Mutable Delegation Chains}
% -------------------------------------------------------

The choice between fixed and mutable delegation chains is not merely an implementation detail --- it determines what accountability guarantees a system can make. A \emph{mutable delegation chain} permits the parent pointer of an agent to be redirected after provisioning: the agent may be re-parented to a different node, its scope of authority may be redefined by an external actor, or its human anchor reference may be updated. Mutable chains enable dynamic task reassignment but break the correspondence between an agent's current parent and its original authorization, making the chain of accountability incomplete or incorrect.

Mutable chains are the norm in conventional multi-agent frameworks. The argument for mutability is operational flexibility: in dynamic environments, task requirements change, agents fail, and organizational structures evolve. A static parent pointer becomes a liability when the parent is decommissioned. The conventional response is governance: policy documents specify conditions under which parent pointers may be redirected, and audit logs record redirection events. This approach works in low-stakes environments but fails in high-stakes ones, where the cost of an unauthorized redirection may be catastrophic and irreversible. The HDP protocol identifies this precisely as the core vulnerability in conventional delegation: without cryptographic chain-of-custody, there is no verifiable link between a terminal action and its original human authorizer \cite{hdp-draft, hdp-draft-2}.

A \emph{fixed delegation chain} --- the approach RSP adopts --- permits no modification to the parent pointer after provisioning. The chain is established once, at spawn time, and persists for the agent's lifetime. This is a stronger constraint than access control policy: the Fact Layer enforces immutability of the parent pointer through write protocol, rejecting any request to overwrite $\alpha(n)$ after provisioning regardless of who issued the request or what justification is offered. The immutability is structural, not contractual.

The fixed chain does not preclude dynamic task reassignment --- it requires that reassignment be implemented through the spawn protocol rather than retroactive pointer modification. When a task must be reassigned, the agent terminates and a new agent is spawned under the appropriate parent. The distinction matters because the spawn event is a Fact Layer write, creating a new provenance record with full metadata, whereas a pointer modification would retroactively obscure the original authorization. The fixed chain preserves the complete history of every authorization decision.

The tradeoff corresponds to what the distributed systems literature terms \emph{strong consistency} versus \emph{eventual consistency} \cite{vogels2009eventual}. Mutable chains are eventually consistent: the current parent pointer reflects the most recent authorized modification, but historical pointers may be overwritten. Fixed chains are strongly consistent: the parent pointer at any time $t$ is identical to the parent pointer at provisioning. RSP chooses strong consistency for the parent pointer because the accountability guarantee requires it. Eventual consistency for chain topology would admit the possibility that, at some point in time, the chain is unverifiable --- and unverifiable chain topology means unverifiable accountability.

% -------------------------------------------------------
\subsection{Hierarchical Task Decomposition in AI}
% -------------------------------------------------------

Hierarchical Task Decomposition (HTD) is the dominant paradigm for managing complexity in AI planning and agentic systems. The core insight, traceable to classical planning research, is that a complex task can be recursively decomposed into simpler subtasks until each leaf is directly executable by a primitive action or a single agent \cite{nilsson1982principles}. Hierarchical Task Network (HTN) planning, the most influential formalization, organizes this decomposition through a task network --- a directed acyclic graph (DAG) of compound and primitive tasks, where edges represent decomposition relationships \cite{ghallab2004htn}. Methods specify how to decompose a compound task into subtasks; conditions specify when a given method applies. The planner selects and applies methods until all tasks are primitive, yielding an executable action sequence.

In classical planning, HTD operates on a task network constructed at plan time, before execution begins, and the network remains fixed during execution. This static structure is well-suited for fully observable, deterministic environments and poorly suited for agentic AI systems, which must operate in partially observable, non-deterministic environments where the appropriate decomposition is discovered during execution.

Modern agentic frameworks have extended HTD to \emph{dynamic decomposition}: agents spawn children at runtime based on conditions encountered during execution, producing a task tree grown rather than pre-constructed. ReAcTree demonstrates that hierarchical LLM agent trees with explicit control flow significantly outperform flat planning approaches --- achieving 61\% goal success on long-horizon benchmarks versus 31\% for flat ReAct --- by decomposing complex goals into subgoal trees managed by specialized agent nodes \cite{reactree2025}. HiPER introduces hierarchical reinforcement learning for LLM agents that decouples planning from execution: a high-level planner proposes subgoals and a low-level executor carries out actions across multiple steps, achieving 97.4\% on ALFWorld and 83.3\% on WebShop by enabling explicit temporal abstraction and hierarchical credit assignment \cite{hiper2025}. The FIRMHIVE framework operationalizes dynamic decomposition as a decentralized Runtime Tree of Agents where agents recursively spawn specialized sub-agents to perform targeted analysis, achieving approximately 16 times more reasoning steps and 5.6 times more security alerts than baseline LLM-agent systems \cite{firmhive2025}. Structured Test-Time Scaling provides a theoretical analysis of why these hierarchical approaches succeed: three-layer structural decoupling (topology compression, scope isolation, strict verification) reduces the effective failure exponent from linear $\Theta(W)$ to near-logarithmic $\tilde{O}(\log W)$, enabling reliable long-horizon reasoning as task complexity grows \cite{structuredttt2026}.

THREAD (Thinking Recursively and Dynamically) provides perhaps the closest conceptual relative to RSP, treating LLM generation as a dynamic multi-threaded process where a parent thread can spawn child threads to handle sub-tasks --- think, retrieve data, plan steps --- and pause until those children return concise results. This dynamic recursive spawning enables hierarchical task decomposition while managing context length, achieving state-of-the-art performance on agent-task and data-grounded QA benchmarks with smaller models \cite{thread2025}. The key insight RSP contributes to this tradition is that decomposition itself is not inherently problematic --- it is the absence of constraints on decomposition that creates risk. Dynamic HTD in open-loop agentic systems has been identified as a vector for \emph{autonomous drift}: the progressive expansion of agent behavior beyond authorized scope, precisely because decomposition is unconstrained \cite{autoGPT2023}. The spawn necessity threshold --- the Purpose Layer's requirement that each child justify why it cannot be replaced by the parent --- addresses a specific HTD failure mode: \emph{artificial fragmentation}, where a parent spawns children not to exploit genuine parallelism or specialized capability but to circumvent context-window limits or create the appearance of multi-agent orchestration. This mirrors \emph{serialization of parallelizable tasks} in classical planning --- the decomposition of a task that should execute concurrently into a sequential chain consuming more depth without delivering more capability. The Purpose Layer's spawn necessity check is a structural deterrent: the justification requirement forces the parent to articulate why spawning is necessary, making artificial fragmentation auditable and rejectable.

% -------------------------------------------------------
\subsection{The BX3 Framework as Substrate for Immutable Parent Chains}
% -------------------------------------------------------

The BX3 Framework's three-layer model --- Purpose Layer, Bounds Engine, and Fact Layer --- provides the minimal sufficient architectural substrate for RSP's accountability guarantees \cite{bx3-framework}. The mapping between RSP components and BX3 layers is not a design choice among alternatives; it is a structural necessity. Each RSP guarantee requires exactly the capabilities that one BX3 layer provides and no other layer can provide.

The \emph{immutable parent pointer} is a Fact Layer write protocol constraint. The Fact Layer is the only layer in the BX3 architecture capable of producing non-reversible physical effects. The parent pointer $\alpha(n)$ is established once, at node provisioning, through a Fact Layer write operation that is atomic and immediately hash-chained to the preceding ledger entry. After this write, no actor --- including the Purpose Layer that originated it --- can modify $\alpha(n)$. This immutability is not a policy enforced by convention; it is a structural constraint enforced by the Fact Layer's write validation. Any attempt to overwrite $\alpha(n)$ is rejected at the protocol level, before execution. This contrasts with ACL-based immutability in conventional systems: administrative compromise in ACL-based systems leads directly to immutability violation, whereas BX3's protocol-level enforcement has no administrative override path \cite{bx3-framework}.

The \emph{depth bound} $D_{\max}$ is a Purpose-Layer-defined, Bounds-Engine-enforced, Fact-Layer-hardened constraint. The Purpose Layer defines the maximum permissible chain depth as part of the spawn authorization policy. The Bounds Engine evaluates each proposed spawn against this limit, returning \texttt{SPAWN\_REFUSED} if the spawn would exceed it. The Fact Layer's pre-commit hook enforces this limit structurally: any spawn operation that would produce $|C| \geq D_{\max}$ is rejected before execution, regardless of the Bounds Engine's return. The constraint is thus enforced three times --- by the layer that defines it, the layer that evaluates against it, and the layer that structurally blocks violations.

The \emph{forensic ledger} is a Fact Layer artifact. Every spawn event, state transition, Purpose Layer directive, and exception condition is recorded in the append-only ledger maintained by the Fact Layer. Each entry is hash-chained to its predecessor, making the ledger tamper-evident: any insertion, deletion, or modification of a ledger entry breaks the hash chain and is detectable by any observer with read access. The ledger provides the evidentional foundation for RSP's correctness properties --- drift prevention, chain integrity, and termination --- because it contains a complete, immutable record of every decision point in the chain's lifecycle.

The \emph{human anchor} $h(n)$ is a Purpose Layer assignment. At provisioning, the Purpose Layer records the identity of the human principal who authorized the root spawn event. This assignment is immutable for the lifetime of the chain: for all nodes $n$ in chain $C$, $h(n) = h(n_0)$ regardless of chain depth. This is the architectural expression of the upstream BX3 accountability guarantee: systems built on the BX3 Framework fail upward into human consciousness, never downward into autonomous chaos.

The structural necessity of the three-layer mapping is established formally in Section \ref{sec:rs-integration}. Removing any layer causes the corresponding guarantee to fail. Without the Purpose Layer, there is no authoritative human anchor and no source of spawn authorization policy. Without the Bounds Engine, depth limits become unenforceable conventions subject to override by operational urgency. Without the Fact Layer, immutability of the parent pointer and tamper-evidence of the ledger are promises rather than constraints. The three-layer architecture is not a design pattern --- it is the minimal structure that makes RSP's correctness properties unconditionally guaranteed rather than conventionally expected.
